\section{Character moments at the two exceptional configurations}
\label{sec:dual}

This section supplies the power saving needed at the two boundary
configurations of the dispersion argument.  The coefficients have a
restricted but important form: they are convolutions of a bounded number
of prime sequences with independently specified smooth norm weights.
There is no sharp condition coupling the prime factors.

Fix a positive constant $c$, a positive integer $k_0$, and fixed support
and smoothness bounds.  An admissible prime convolution at length $Y$ is
\begin{equation}
 \beta(b)=\mu^2(b)
 \sum_{b=\pi_1\cdots\pi_k}
       \prod_{j=1}^k w_j\bigl(N(\pi_j)/F_j\bigr),
 \qquad 1\le k\le k_0,\qquad \prod_{j=1}^k F_j\asymp Y,
 \label{eq:prime-convolution}
\end{equation}
where the primes are primary, the tuple is ordered, every supported prime
has norm at least $Y^c$, and the $w_j$ have compact support in fixed
subintervals of $(0,\infty)$.  Each fixed derivative is bounded by a fixed
power of $\log Y$, uniformly over the family under consideration.  These
bounds imply fixed divisor bounds for all coefficients below.

For sums $\sum_b\beta(b)\chi_b(pq^2)$, with $p,q$ coprime and
squarefree of norm lengths $P,Q$, the two sieves reach the same bound
at different configurations. At $(P,Q)=(Y,1)$,
Theorem~\ref{thm:cls} gives, up to an arbitrarily small power loss,
the second-moment bound
\[
 Y\bigl(Y+Y+Y^{4/3}\bigr)\ll Y^{7/3}.
\]
At $P=Q=Y^{1/3}$ the conductor norm is $O(PQ)=O(Y^{2/3})$, so
Lemma~\ref{lem:ordinary-sieve} gives $Y(Y+Y^{4/3})\ll Y^{7/3}$,
again up to an arbitrarily small power loss. These are the two
configurations left by the dispersion argument below. The convolution
structure permits a fixed power saving in neighborhoods of both.

\begin{proposition}[The exceptional character moments]
\label{prop:dual}
There are constants $\delta,\eta>0$, depending only on the fixed data in
\eqref{eq:prime-convolution}, with the following property.  Let
$v,e\in\mathcal O\setminus\{0\}$ satisfy
$N(v),N(e)\le Y^\delta$, and let $1+|u|\le Y^{0.36}$.  Put
\[
 A(p,q)=\sum_{(b,e)=1}\beta(b)N(b)^{iu}\chi_b(vpq^2).
\]
If $p,q$ run through coprime primary squarefree elements of respective
norm lengths $P,Q\ge1$, and either
\[
 \left|\frac{\log P}{\log Y}-1\right|\le\delta,
 \qquad 0\le\frac{\log Q}{\log Y}\le\delta,
\]
or
\[
 \left|\frac{\log P}{\log Y}-\frac13\right|\le\delta,
 \qquad
 \left|\frac{\log Q}{\log Y}-\frac13\right|\le\delta,
\]
then
\begin{equation}
 \sum_{p,q}|A(p,q)|^2\ll Y^{7/3-\eta}.
 \label{eq:exceptional-moment}
\end{equation}
The implied constant is uniform for the specified support and
smoothness bounds.
\end{proposition}

The short convolution identity used in the proof replaces each prime
sum by products of truncated Mobius sums and full smooth sums.
A factor of essentially full length is shortened by the functional
equation. Otherwise, moments of products with factors repeated or
omitted rule out the second configuration and force a precise relation
between factor length and value in the first. There the off-diagonal
Gram estimate supplies the final contradiction. We first prove the two
analytic estimates needed for these alternatives. Throughout, the
unrestricted-inner form of Lemma~\ref{lem:cubic-extensions} permits
divisor-bounded coefficients that need not be squarefree.

In the next lemma a full sum means the complete ideal sum of the
primitive Hecke $L$-function. A restriction excluding primes above $3$,
or other missing Euler factors, is first removed by inclusion--exclusion;
this is carried out explicitly in the dominant-full-sum argument below.

\begin{lemma}[A full smooth sum and its dual]
\label{lem:smooth-character-duality}
Let $\chi$ be a primitive nonprincipal Hecke character of trivial
infinite parameter and conductor norm $D$. Let $W$ be smooth and supported in a
fixed compact subinterval of $(0,\infty)$, with each fixed derivative
bounded by a fixed power of $\log Y$. The full sum
$\sum_n W(N(n)/Z)\chi(n)N(n)^{it}$ admits a decomposition into dual
norm dyads
\[
 J\ll Y^\epsilon D(1+|t|)^2/Z
\]
with amplitude $O(Y^\epsilon\sqrt{Z/J})$.  The coefficients in a dual
dyad are divisor-bounded, and its character is $\overline\chi$.
All additional norm shifts can be integrated against a common
conductor-independent $L^1$ majorant of mass $O(Y^\epsilon)$.
The omitted tail is smaller than any prescribed negative power of $Y$,
provided the lengths and heights are bounded by fixed powers of $Y$.
\end{lemma}

\begin{proof}
Write
$\widetilde W(s)=\int_0^\infty W(x)x^s\,dx/x$.  Fixed logarithmic
factors in the coefficients can be included in $W$.
Mellin inversion expresses the sum as
\[
 \frac1{2\pi i}\int \widetilde W(s)Z^s L(s-it,\chi)\,ds.
\]
The character is nonprincipal, so shifting the contour to the left
crosses no pole of the $L$-function.  Up to fixed field constants, its
functional equation is
\[
 L(s,\chi)=\epsilon(\chi)D^{1/2-s}
       \frac{\Gamma(1-s)}{\Gamma(s)}L(1-s,\overline\chi),
 \qquad |\epsilon(\chi)|=1.
\]
On a sufficiently far-left line the dual Dirichlet series is absolutely
convergent.  Partition it into smooth norm dyads there.  Stirling's
formula and the rapid decay of $\widetilde W$ show that dyads beyond
$Y^\epsilon D(1+|t|)^2/Z$ have an arbitrarily small total tail.
This truncation is performed before moving the retained, finite dyads
back to the line $\Re s=1/2$.

On that line, with $s=1/2+i\tau$, the integrand for a retained dyad is
$\sqrt Z\,\widetilde W(1/2+i\tau)$ times
\begin{equation}
 \epsilon(\chi)D^{i(t-\tau)}Z^{i\tau}
 \frac{\Gamma(1/2-i(\tau-t))}
      {\Gamma(1/2+i(\tau-t))}
 \sum_{N(n)\asymp J}
       \overline\chi(n)N(n)^{-1/2+i(\tau-t)}w_J(N(n)/J).
 \label{eq:symmetry-line-character}
\end{equation}
The multiplier preceding the polynomial has absolute value one.
In particular, its conductor and root-number dependence is a scalar
for each character, not a change in the coefficients of the polynomial.
The gamma poles lie to the right of this move.  Extracting $J^{-1/2}$
from the polynomial proves the amplitude assertion.  Rapid Mellin
decay gives a common $L^1$ majorant in $\tau$, including any fixed
logarithmic derivative costs.  Minkowski's inequality in a square-sum
norm over characters therefore applies with the same shifts for every
character.  A common cutoff based on the largest conductor in the
average is permissible by the earlier far-left tail estimate.
\end{proof}

\begin{lemma}[An off-diagonal Gram estimate]
\label{lem:character-gram}
Let $p$ be primary squarefree with $N(p)\asymp P$, and let $W$ be a
fixed smooth radial norm weight of compact support in $(0,\infty)$.
For primary squarefree $p'$ of the same norm length put
\[
 T_{p,p'}=\sum_{n\equiv1\, (3)}
     W(N(n)/Z)\chi_p(n)\overline{\chi_{p'}(n)}.
\]
For $Z\ge1$ and every $\epsilon>0$,
\begin{equation}
 \sum_{p'\ne p}|T_{p,p'}|^2
 \ll_\epsilon (PZ)^\epsilon
      Z\left(P+(P^3/Z)^{2/3}\right).
 \label{eq:near-orthogonality}
\end{equation}
The same bound holds with a subset of the available $p'$.
\end{lemma}

\begin{proof}
Fix $k=(p,p')$, put $K=N(k)$, and write $p=ka$, $p'=kb$.
The remaining moduli $a,b$ are coprime and squarefree.  The common
factor contributes the restriction $(n,k)=1$.  Remove this restriction
by summing over $m\mid k$ with coefficient $\mu(m)$, and put $M=N(m)$.
After $n=mx$, the character is
$\psi(x)=\chi_a(x)\overline{\chi_b(x)}$, apart from a scalar of
absolute value one.  It is primitive modulo $ab$ and nonprincipal:
the only principal case would be $a=b=1$, contrary to $p'\ne p$.

Apply Poisson summation on $x\equiv1\pmod3$ modulo $3ab$.
The zero frequency vanishes.  Primitive Gauss evaluation gives
amplitude
\[
 \asymp\frac{Z}{M\sqrt{N(ab)}}\asymp\frac{ZK}{MP},
\]
and natural dual norm length
\[
 J_0\asymp\frac{N(ab)M}{Z}
       \asymp\frac{P^2M}{K^2Z}.
\]
For clarity about the subsequent sieve application, the complete
Gauss factor and CRT unit factors are independent of the frequency.
They are scalars for each $b$.  The frequency character is
$\overline{\chi_a(h)}\chi_b(h)$, and the primary residue condition
adds only a fixed bounded-modulus phase.  The Fourier transform is
radial, so its remaining dependence on $b$ is through its norm.

More explicitly, set $c=ab$ and
$g(\psi)=\sum_{x\bmod c}\psi(x)e(\operatorname{Tr}(x/c))$.
Writing the primary residue classes as $x=c+3s$ gives the complete
finite Fourier coefficient
\[
 e\left(\operatorname{Tr}\frac{h}{3\lambda}\right)
       \psi(3\lambda)\overline{\psi(h)}g(\psi)
 =e\left(\operatorname{Tr}\frac{h}{3\lambda}\right)
       \overline{\psi(h)}g(\psi),
\]
because $3\lambda=(-\lambda)^3$. This also displays directly why
no varying-modulus factor remains inside the frequency coefficients.

On a frequency dyad $N(h)\asymp J$, Mellin separation in the compact
scaled variables $N(h)/J$ and $N(b)/(P/K)$ gives coefficients
independent of $b$, integrated against a rapidly decreasing common
majorant.  For every fixed $A>0$ its mass is
$O_A((1+J/J_0)^{-A})$; derivatives have the same property after
adjusting $A$.  Units and exact powers of the ramified prime in $h$
can be separated first.  Their character values are again rowwise
scalars, and the number of cases costs an arbitrarily small power.
The extension with unrestricted inner coefficients in
Lemma~\ref{lem:cubic-extensions} now bounds this dyad, after squaring
and summing over $b$, by
\[
 (PZ)^\epsilon
 \left(\frac{ZK}{MP}\right)^2
 J\left(\frac PK+J+(PJ/K)^{2/3}\right)
 (1+J/J_0)^{-A}.
\]
The restrictions on $b$ may be dropped in this upper bound.

Sum the frequency dyads, retaining the rapid decay also when
$J_0<1$.  The three powers of $J$ are $1,2,5/3$; for $A$ sufficiently
large their dyadic sums are bounded by the corresponding powers of
$J_0$.  Substitution gives the three terms
\[
 \frac{ZP}{MK},\qquad \frac{P^2}{K^2},\qquad
 \frac{P^2Z^{1/3}}{K^2M^{1/3}}.
\]
Since $Z,K,M\ge1$, their sum is
$O(ZP+P^2Z^{1/3})$.  There are only a divisor-bounded number of
choices $k\mid p$ and $m\mid k$.  Cauchy's inequality over those
choices, with a rechoice of $\epsilon$, proves
\eqref{eq:near-orthogonality}.
\end{proof}

\begin{proof}[Proof of Proposition~\ref{prop:dual}]
We use a limiting-exponent contradiction, and give the uniformity
argument explicitly at the end.  Thus suppose that along a sequence
$Y\to\infty$ the extra twist and exclusion norms are $Y^{o(1)}$,
the pair of length exponents tends to $(1,0)$ or $(1/3,1/3)$, and
the moment is at least $Y^{7/3-o(1)}$.  Throughout this proof, a
subpower loss means a bound with an arbitrarily small fixed positive
power loss; it is never used as a logarithmic saving.

\smallskip
\noindent\emph{Prime powers and repeated factors.}
First remove the squarefree restriction in
\eqref{eq:prime-convolution}, and replace each prime weight by its
smooth $\Lambda/\log N$ counterpart.  The resulting discrepancy has
divisor-bounded coefficients on a set of size $O(Y^{1-\theta})$ for
some fixed $\theta>0$.  Indeed every discrepancy contains either a
repeated prime of norm at least $Y^c$, or a prime power of exponent
at least two and norm at least $Y^c$.  In either case its squareful
part has positive-power norm.  In the canonical factorization
$b=st^2d^3$, with $s,t$ coprime squarefree, this also implies
$N(s)\le Y^{1-\theta'}$ for some fixed $\theta'>0$.
To include prime powers with a small base, if $\pi^j$ has $j\ge2$
and $N(\pi^j)\ge Y^c$, then
$N(b/s)\ge N(\pi^j)^{1/2}\ge Y^{c/2}$.
Also $b$ is divisible by the square of
$\pi^{\lfloor j/2\rfloor}$, whose norm is at least $Y^{c/4}$.
Summing $O(Y/N(d)^2)$ over square divisors with $N(d)\ge Y^{c/4}$
gives support size $O(Y^{1-c/4})$. Repeated large primes satisfy the
same bounds, with room to spare.

In the first configuration fix $q$, at a subpower cost. In the
factorization $b=st^2d^3$, let $S,V,C$ be the dyadic norms of $s,t,d$.
Then $SV^2C^3\asymp Y$ gives $(VC)^2S\ll Y$, while
$S\ll Y^{1-\theta'}$. Substituting these
two bounds in \eqref{eq:inner-extension-dyad}, and summing the dyads,
bounds the discrepancy moment by
\[
 Y^{1+o(1)}
 \left(Y+Y^{1-\theta'}+Y^{(4-2\theta')/3}\right),
\]
which saves a power.  In the second configuration, square the
discrepancy polynomial.  Its support has size at most
$Y^{2-2\theta}$ and its coefficients remain divisor-bounded.
The ordinary character sieve, at polynomial length $Y^2$ and
conductor norm at most $Y^{2/3+o(1)}$, therefore bounds its fourth
moment by $Y^{4-2\theta+o(1)}$.  Cauchy's inequality over
$Y^{2/3+o(1)}$ characters gives a second moment at most
$Y^{7/3-\theta+o(1)}$.  The fixed extra character and exclusions are
part of the coefficients in these applications.  Hence both
discrepancies can be discarded.

\smallskip
\noindent\emph{A short convolution identity.}
Here $*$ denotes Dirichlet convolution on ideals and $1$ is the
constant-one arithmetic function. For a prime-polynomial dyad with
upper endpoint $F'$, let
$m=\mu\,1_{N\le\sqrt{F'}}$ and let $\mathbf1_*$ denote the
convolution identity.  The function $\mathbf1_*-m*1$ vanishes up to
norm $\sqrt{F'}$, so its convolution square vanishes up to norm
$F'$.  Convolving with $\Lambda$ and using $1*\Lambda=\log N$
gives, on the dyad,
\begin{equation}
 \Lambda=(2m-m*m*1)*\log N.
 \label{eq:short-mobius-identity}
\end{equation}
Apply this identity to every factor, insert smooth norm dyads, and
separate the original product weights by Mellin inversion.  The
number of pieces and the $L^1$ costs are subpower.  Rapid Mellin
decay permits truncation of additional shifts at $Y^\epsilon$;
the tails are negligible by trivial bounds.  Minkowski's inequality
permits fixing the same shifts across each character average.

It follows that some product of a bounded number of divisor-bounded
character polynomials still has second moment
$Y^{7/3-o(1)}$.  Pass to a subsequence on which their length
exponents tend to
\[
 a_i\ge0,\qquad \sum_i a_i=1.
\]
The common norm twists and fixed character factors can be distributed
multiplicatively over these polynomials.  Every truncated Mobius
factor has length at most $Y^{1/2+o(1)}$.  Thus, if an $a_i$ is one,
that factor is a full smooth sum, with logarithmic factors allowed,
and with twist height at most $Y^{0.37}$.

\smallskip
\noindent\emph{A dominant full sum.}
Suppose that $a_i=1$.  All other factors have subpower trivial size.
The character in the full sum is nonprincipal.  Its conductor away
from the fixed ramified part agrees with $pq$ except at primes
dividing the subpower extra twist or exclusion number.  Partition
according to the part of $pq$ at those primes and their local
character values, and remove missing Euler factors by
inclusion-exclusion.  There are subpowerly many possibilities.
For each one the rescaled full-sum length is still $Y^{1+o(1)}$;
outside the fixed small parts, the variable character is the original
cubic character or its conjugate.  Exact powers of the ramified
prime can likewise be separated on the dual side.
In detail, let $S$ be the primes dividing $3ve$, fix the $S$ parts
of $p,q$, and fix the resulting primitive local character $\psi_S$.
There are $Y^{o(1)}$ choices, since the product of the relevant small
prime norms is subpower. The remaining primitive character is
$\psi_S\chi_{p_0}\chi_{q_0}^2$, with $(p_0q_0,S)=1$.
Inclusion-exclusion for missing Euler factors rescales by a fixed
subpower divisor before the functional equation is applied.
On the dual side $\overline{\psi_S}$ is a fixed coefficient twist;
its root number and conductor phase are rowwise scalars as in
\eqref{eq:symmetry-line-character}. The surviving $p_0$, or $p_0q_0$,
has positive-power norm, so this primitive character is nonprincipal.

Apply Lemma~\ref{lem:smooth-character-duality}.  In the second
configuration the conductor is at most $Y^{2/3+o(1)}$, so the trivial
dual estimate is
\[
 |\text{full sum}|\ll Y^{1/3+0.37+o(1)}.
\]
Summing squares over the available characters has exponent at most
$2/3+2(1/3+0.37)<7/3$.  In the first configuration fix $q$ and the
small parts.  The variable squarefree conductor has length
$Y^{1+o(1)}$, and the dual polynomial has length
$J\le Y^{0.74+o(1)}$, allowing an arbitrarily small fixed power
loss.  The cubic sieve with unrestricted inner coefficients gives
\[
 Y^{1+o(1)}\left(Y+J+(YJ)^{2/3}\right)
 \le Y^{2.16+o(1)}.
\]
Both estimates contradict the unsaved moment.  The averaging here is
legitimate precisely because the conductor-dependent factors in
\eqref{eq:symmetry-line-character} are outside each polynomial.

\smallskip
\noindent\emph{Moments of different lengths.}
We may now assume that every $a_i<1$.  Pigeonhole the absolute values
of the individual polynomials.  Values smaller than a sufficiently
large fixed negative power of $Y$ can be discarded using the trivial
bounds for the other factors and the number of characters.  After a
further subsequence, there is a set of $R$ rows on which their
absolute values are $Y^{v_i+o(1)}$, where $v_i\le a_i$, and
\begin{equation}
 r:=\lim\frac{\log R}{\log Y}\ge\frac73-2v,
 \qquad v=\sum_i v_i.
 \label{eq:dual-large-values}
\end{equation}

In the second configuration, cardinality gives $r\le2/3$.
The ordinary sieve applied to the square of the full product gives
$r+4v\le4$.  Together with \eqref{eq:dual-large-values}, these imply
$v=5/6$ and $r=2/3$.  There is therefore an index with
$0<a_i<1$ and $v_i\ge5a_i/6$.  Multiply the full product by this
additional factor, and put $z=1+a_i\in(1,2)$.  Its value on the
chosen rows has exponent at least $5z/6$.  The ordinary sieve yields
\[
 r\le \max(4/3,z)+z-5z/3
      =\max(4/3,z)-2z/3<2/3,
\]
a contradiction.

In the first configuration fix $q$ without changing the limiting
inequality \eqref{eq:dual-large-values}.  The rows are now distinct
squarefree $p$ of length $Y^{1+o(1)}$.  Put
$d_i=v_i-5a_i/6$.  For any fixed nonnegative integer multiplicities
$k_i$ such that $z=\sum_i k_i a_i\in[1/2,2]$, the corresponding
product is a divisor-bounded polynomial of length $Y^{z+o(1)}$.
The unrestricted-inner cubic sieve gives
\begin{equation}
 r\le \frac23-2\sum_i k_i d_i.
 \label{eq:dual-multiplicity}
\end{equation}
Here we used
\[
 z+\max\{1,z,2(1+z)/3\}=\frac23+\frac{5z}{3}
 \qquad (1/2\le z\le2).
\]
The all-ones multiplicity vector makes the right side of
\eqref{eq:dual-multiplicity} equal to the lower bound in
\eqref{eq:dual-large-values}.  Increasing any one multiplicity
therefore proves $d_i\le0$.  Decreasing it proves $d_i\ge0$ whenever
$a_i\le1/2$.  Thus all short factors have $d_i=0$, including factors
with $a_i=0$.

There is at most one index with $a_i>1/2$.  If there is one, remove
it and use copies of a positive short factor to obtain a length
exponent in $[1/2,2]$.  Such a short factor exists because no $a_i$
is one.  Equation~\eqref{eq:dual-multiplicity} then gives
$r\le2/3$, forcing the remaining $d_i\ge0$ as well.  Consequently
all $d_i=0$ and $r=2/3$.

Choose a positive short exponent $a\le1/2$.  Some fixed integer
multiple of $a$ lies strictly between $4/3$ and $2$, since the
width of that interval exceeds $a$.  Taking the corresponding
copies produces a divisor-bounded polynomial of length
$Z=Y^{z+o(1)}$, with $4/3<z<2$, whose values are
$Y^{5z/6+o(1)}$ on at least $Y^{2/3-o(1)}$ distinct squarefree rows.
The integer multiplicity is fixed after passing to the subsequence;
it does not grow with $Y$, even when the limiting $a$ is small.

\smallskip
\noindent\emph{Near orthogonality.}
Choose a nonnegative smooth radial weight equal to at least one on
the support of this last polynomial.  Weighted Cauchy and the
row-sum bound for its Gram matrix give the classical
Bombieri--Hal\'asz--Montgomery inequality; see
\cite[Theorem~1, equation~(1)]{HarcosBHM} for this Hilbert-space step.
In the present notation it gives
\[
 RY^{5z/3+o(1)}
 \le ZY^{o(1)}
       \left(O(Z)+\max_p\sum_{p'\ne p}|T_{p,p'}|\right).
\]
The support ratio is fixed, since the number of copied factors is
fixed.  Lemma~\ref{lem:character-gram}, with
$P=Y^{1+o(1)}$, bounds each squared off-diagonal row sum by
\[
 Y^{o(1)}Z\left(Y+(Y^3/Z)^{2/3}\right).
\]
After dividing the Gram inequality by $ZR$, the left side has
exponent $2z/3$.  The diagonal has exponent at most $z-2/3$, and
Cauchy's inequality bounds the off-diagonal average by exponent
\[
 \frac{z+\max(1,2-2z/3)-2/3}{2}.
\]
Both are strictly less than $2z/3$: the diagonal gap is $(2-z)/3$;
the other gap is $(3z-4)/6$ for $z\le3/2$, and $(z-1)/6$ for
$z\ge3/2$.  These positive gaps absorb every arbitrarily small
fixed power loss and give the final contradiction.

\smallskip
\noindent\emph{Uniform neighborhoods.}
If the proposition failed for every positive neighborhood size and
power saving, choose successively shrinking neighborhood sizes and
savings, and then arbitrarily large counterexamples.  A diagonal
sequence would have extra twist and exclusion norms $Y^{o(1)}$,
length exponents tending to one of the specified pairs, and moment
at least $Y^{7/3-o(1)}$.  There are only boundedly many factors, and
their length exponents lie in a compact simplex.  The preceding
pigeonholing also confines all relevant value exponents to a fixed
compact interval.  The proof therefore applies to a subsequence of
these counterexamples.  Every selected multiplicity and every
small power loss is fixed before taking the limit.  The resulting
contradiction proves the existence of fixed $\delta,\eta>0$ and
the asserted uniform bound.
\end{proof}
